On applique une tension $t\mapsto 5+\e^{-t}$ aux bornes d'un circuit. Le courant vérifie l'équation différentielle : $y'+5y=5+\e^{-t}$

\begin{enumerate}
	\item Donner et résoudre l'équation homogène (E$_0$)\par \vspace*{0.5cm}

	\Lignes{2}
	\item A l’aide du logiciel Géogébra, déterminer les nombres réels $a$ et $b$ pour que la fonction $p$ définie sur $\R$ par : $p(x) = a + b~\e^{-t}$ soit une solution particulière de l'équation différentielle (E), et donner l’expression de $p$. 
	
	On trouve : $a=$\makebox[0.1\linewidth]{\dotfill} \hfill $b=$\makebox[0.1\linewidth]{\dotfill} \hfill $p(x)=$\makebox[0.3\linewidth]{\dotfill}

	\smallskip
	
	\appelprof
	\item En déduire la forme générale des solutions de (E).\par \vspace*{0.5cm}

	\Lignes{1}
\end{enumerate}